\begin{table}[t]
\centering
\caption{Baselines are shown at the tightest realised false-alarm probability they achieve among settings with a miss rate below 0.9, that is, detecting at least one event in ten (the miss-rate column gives the value achieved). The two exact e-processes are shown at a common nominal level $\alpha=0.01$, so they are compared at equal nominal level. Streams are 60 minutes long (200 null and 200 post-change streams per operating point). This is the operator's question: insisting on few false alarms, what does the method cost in delay? Both exact e-processes tested (SENTINEL-E and the changepoint mixture it strictly contains, \cref{prop:generalisation}) reach a rate of at most 0.05; none of the classical or parametric baselines does. The classical procedures can be made fast, but only at a false-alarm probability near one for most settings of their knob; the full trade-off is in Figure~\ref{fig:delayfar}.}
\label{tab:delay_far}
\resizebox{\ifdim\width>\linewidth\linewidth\else\width\fi}{!}{%
\begin{tabular}{llrcrrrr}
\toprule
method & nominal $\alpha$ / knob & realised PFA & $\le 0.05$? & ARL$_0$ & cond. delay (s) & censored delay (s) & miss rate \\
\midrule
SENTINEL-E & 0.01 & 0.000 & \checkmark & 90,000 & 205.6 & 244.5 $\pm$ 65.2 & 0.015 \\
Fixed threshold & $10^{-5}$ & \textbf{0.105} & \ding{55} & 85,085 & 1489.3 & 2649.3 $\pm$ 71.7 & 0.885 \\
p-value threshold & 0.003 & \textbf{0.680} & \ding{55} & 53,943 & 297.9 & 598.1 $\pm$ 122.8 & 0.120 \\
CUSUM & 20 & \textbf{0.990} & \ding{55} & 16,982 & 389.1 & 413.2 $\pm$ 76.3 & 0.010 \\
Shiryaev--Roberts & 10 & \textbf{1.000} & \ding{55} & 619 & 4.3 & 4.3 $\pm$ 0.4 & 0.000 \\
Parametric e-detector & 0.5 & \textbf{1.000} & \ding{55} & 4,444 & 17.4 & 17.4 $\pm$ 7.0 & 0.000 \\
E-SHIFT & 0.002 & \textbf{0.965} & \ding{55} & 12,136 & 279.3 & 1791.7 $\pm$ 178.6 & 0.600 \\
Changepoint mixture & 0.01 & 0.000 & \checkmark & 90,000 & 166.5 & 311.3 $\pm$ 92.5 & 0.055 \\
\bottomrule
\end{tabular}}
\begin{minipage}{\linewidth}\vspace{2pt}\footnotesize Delays are in seconds at 25\,fps. For SENTINEL-E and the changepoint mixture the knob is the nominal level $\alpha$; for the baselines it is their own threshold, per-frame rate or betting level. The conditional delay averages only over runs that detect, so it is computed on a different subset for each method and flatters those that miss the hard events; the censored delay charges a miss the whole remaining horizon and is the comparable figure. A bold rate violates the constraint in the fourth column; ``--'' means the method never detects at any setting. The $\pm$ values are 1.96 standard errors of the censored-delay mean over streams.\end{minipage}
\end{table}
